\section{Data}
\label{sec:data}

\paragraph{Universe.}
The universe is the current list of S\&P~500 constituents, taken from Wikipedia and saved as a
snapshot file so that every run uses the same list. Of the 503 listed share classes, Yahoo
Finance returns data for \NStocksTotal{}. The three missing are recent spin-offs with no
history in the sample.

\paragraph{Prices and returns.}
Daily closing prices, adjusted for splits and dividends, and daily share volumes are downloaded
from Yahoo Finance for \DataStart{} to \DataEnd{}. The return of stock $i$ on day $t$ is the
simple total return
\begin{equation}
  R_{it} = \frac{P_{it}}{P_{i,t-1}} - 1 .
\end{equation}
Missing prices are \emph{not} filled. A gap leaves a missing return, and the stock cannot be
traded until it has a complete estimation window again (\cref{sec:method}). Yahoo Finance
backfills prices for some companies from before their US listing (e.g.\ SW, FERG, AMCR). These
days carry zero reported volume, so we treat every zero-volume day as a non-trading day.

\paragraph{Survivorship bias.}
Using today's index members for a historical backtest introduces survivorship bias. Stocks that
were removed from the index, acquired or delisted between 2010 and today are missing, and
today's members are by construction companies that survived and mostly grew. For a
market-neutral strategy the direction of the bias is not obvious, since it affects both long and
short legs. It does, however, make the universe less representative and removes exactly the
stocks with the most extreme idiosyncratic moves. All results in this draft should be read with
this caveat. The data layer is built so that the CRSP daily stock file, which provides a
point-in-time universe and delisting returns, can replace Yahoo Finance without changes to the
strategy code (\cref{sec:impl}).
\todo{Re-run on CRSP via WRDS and report the difference.}

\paragraph{Sample split.}
The data are split into an in-sample period, which ends on 31~December~2018, and a held-out test
period from \SplitDate{} to \DataEnd{}. Because the first 252 days are needed to estimate the
model, the first trading day is \InSampleStart{}, giving \InSampleDays{} in-sample trading
days. The split and how it is enforced are described in \cref{sec:protocol}.
